\chapter{水波动力学}
本章应用理想\$\textasciitilde{}体动力学理论研究重力场中水面波动现象.

\section*{6.1  水波动力学的基本方程和边界条件}\addcontentsline{toc}{section}{6.1 水波动力学的基本方程和边界条件}

水波是常见的自然现象，如船舰航行时的船首波，风生波，月球引力场周期性变 化产生的潮沙波，地震诱发的海啸，以及盛液容器不规则运动时的液面晃动等.虽然 液面波动千姿百态，它们的波长可以小到几毫米、大到上百米，但是它们都属于不可 压缩流体在重力场中的运动.水波动力学的研究对于海洋、水利等工程和自然环境的 预报是很有用的. 为了便于分析而又抓住主要矛盾，假定水波是不可压缩理想流体在重力场中从 静止状态由外界激励生成.根据凯尔文定理，不可压缩理想流体在重力场(有势力场) 中由静止状态开始的运动是无旋的，就是说理想水波运动属于不可压缩流体的元旋 流动. 研究水波现象，除了要了解流场状态外，更有兴趣的是掌握水面演化的规律，它 是工程计算和自然环境预报中很重要的物理量.通常水波表面高度以原静止水平面 为基准，写成 (x ， y ， z 是直角坐标 ， z 是重力加速度的相反方向，即垂直向上 h

\begin{equation*}z = ~(x ,  y ,  t)\tag{6.1}\end{equation*}

\textasciitilde{}(x ， y ， t) 称为波高画数，简称波高，液面以上是大气，故液面又称为自由面.

综上所述，水波动力学基于以下假定 z (1)水是不可压缩流体，即 v . U=O; (2) 水是无粘理想流体，即水波中的流体应力状 态是 : T;j = - p8;j ; (3) 水波在重力场中运动，重力场的位势为 z z

\[/F\]

第 6 章 水波动力学

\begin{gather*}
n=gz;\\
/ .' /' ~.- ./ /' / / / / / /'
\end{gather*}

(4) 水波在重力场中由静止启动，是无旋的，有速 度势 $\phi$ ，且 U=v$\phi$. 图 6. 1 波面坐标系 根据以上假定，水波动力学属于不可压缩理想流体的非定常无旋运动，它的基本 方程和边界条件如下. 1.基本方程

\[A\phi =0\]

3$\phi$1....\textasciitilde{} ....\textasciitilde{}， ${\rm þ}$ $\tau$::+$\tau$V$\phi$. v$\phi$+ .1二 + gz = C(t) (6.2) (6.3) $\phi$ (x.y.z.t) 是水波流场的速度势，式 (6. 2) 是不可压缩元旋流动的速度势方程， 式 (6.3) 是柯西-拉格朗日积分.式 (6.3) 中积分常数可以用变量替换消去.令 $\phi$'= $\phi$ +fC 川，则速度场 v${\rm ø}$=v${\rm ø}$' 并不改变，这时方程 (6.3) 中空间积分常数 C(t) =0 2. 水波运动的边界条件 (1)固壁不可穿透条件 如水波中存在运动固壁，则在固壁面上应有 (去 t" = 在水波中的静止固壁上 z (2) 自由面(水面〉运动学条件 (翌L.. = $\theta$ n 1 :1:" (6.4a) (6.4b) 水波问题的主要特点是存在运动水面.水面是由水质点组成的流体面，水面上的 质点始终位于水面上，因此水面方程应满足第 2 章的流体面守恒方程，将水面方程 z= t"( x. y , t) 写成

\[4W: F(x.y.z.t) = t"(x.y.t) - z = 0\]

代入流体面守恒方程 (2.13): ${\rm e}$J F/${\rm e}$Jt+U ${\rm •}$ VF=O. 得 \textasciitilde{}十 uEE+UEE-ztJ=0${\rm e}$J t . -- ${\rm e}$Jx . - ${\rm e}$J y (6.5)

式中 (u.v.w) 是水面上质点速度，注意到 V$\phi$ =u. 则上式又可写成 \textasciitilde{}+翌\textasciitilde{}+丝丝一丝 =0dt dx dx dy dy dz 或 (6.6)

\[~ = w (6.7) Dt\]

其中些 =u 坐十 U 坐+当是披高的质点导数，也就是波高垂直位移速度，因此Dt - iJx . - iJ y . iJt 式 (6. 7) 的运动学意义是波面的垂直位移速度等于当地流体质点速度的垂直 分量. (3) 自由面上的压强条件 水波问题需要联立求解速度势方程和柯西-拉格朗日积分，因此必须给定水面 $\tilde{w}$ 上压强边界条件.设水面上方的大气压强为丸，则由于液面的表面张力，自由面下 侧压强 PU' 和大气压强之差为 P. - Pw = r[ 是+ $\tilde{J}$ (6.8) 自由面曲率由波面方程 z= ${\rm C}$(x.y. t) 求出 z 主主 主Z 1 \_ iJz2 1 \_ iJy2

\[R - i1+(~)2 +(~)21 3!2' Ry i1+(~)2 +(~)2r/2\]

L 飞 ax I 飞 ay I J L 飞 ax I 飞 ay I J 当波高 ${\rm C}$max 和波长 $\lambda$ 之比为小量时，即 I iJ ${\rm c}$;iJ x I <<1. I iJ ${\rm c}$;iJ y I <<1. 以及 ${\rm A}$la 2 ${\rm c}$; iJ x 2

\[1<<\]

1.${\rm A}$ I iJ2 ${\rm c}$;iJy 2 1<<1. 则自由面曲率很小，表面张力影响可以忽略.本书只讨论不计表面 张力的水波，这时波面压强条件为

\[z = C(x.y.t). \rho =\rho a\]

将它代入柯西-拉格朗日积分，在自由面上，应有 8$\phi$l 由 $\tau$丁 +$\tau$V$\phi$. v$\phi$+ 旦 +g${\rm c}$ = 0 Ol G P 现在把水波动力学基本方程和边界条件归纳如下. 在水波流场中速度势和压强满足以下方程 z

\begin{gather*}
A\phi =0\\
3\phi 1...~ ...~,  \theta  .~~+\tau V\phi . v\phi  +L+ gz = 0
\end{gather*}

al G P 在固壁边界 $\tilde{B}$ 和水面边界 $\tilde{w}$ 上速度势满足以下边界条件 z $\tilde{R}$: \textasciitilde{}$\phi$ 盯B: 五百一 Ll b - " $\tilde{w}$: 笠+坐坐+坐坐=壁，壁十 i w-w+ 阜 +g${\rm c}$ = 0 ${\rm •}$ iJt . a.r dX . iJ y iJy iJ z' at' 2 . - . - . p (6. 9) (6.10)

水面边界上的常数压强可以用变量替换消去，设 $\phi$=$\phi$，'+$\rho$ ${\rm •}$ t ， 则有 V$\phi$'=v$\phi$ ，变量替 换不影响速度场，这时水面压强条件为

\[3\phi l\]

37+ 言 V$\phi$. v$\phi$+g' = 0 域内柯西-拉格朗日积分为

\[3\phi 1 '(7.... ~ '(7.... I P - P. I -\]

"-\textasciitilde{}+丁 V$\phi$. v$\phi$+ 一一\textasciitilde{}+J!Z=O dt Z P 水波动力学问题的未知量有速度势 $\phi$ 、压强 $\rho$ 和自由面披高 '(x ， y ， 仆，由于存 在自由面，水波动力学流场的解法和无界不可压缩无旋流场的解法有很大差别.第 一，求解速度势的边界条件中含有未知量 '(.r ， y ， t)( 波高函数) ;第二，水波动力学问 题是非线性的，非线性效应主要来自水面的运动学和压强条件.由于以上两个特点， 一般的水波动力学问题的求解较之无界无旋绕流问题要困难一些.

\section*{6.2  等深度水域中小振幅波的线性近似}\addcontentsline{toc}{section}{6.2 等深度水域中小振幅波的线性近似}

现在来讨论一种最简单的水波现象.假定原静止水域是等深度的，即水底方程为

\[Z/. =- h = const. (6. 11)\]

此外，披陡很小，即液面的斜率很小 z |手 1- 问 1- O(e) , e<< 1 (6.12) Jxl l iJy 为了简化分析，还假定水波流场是二维的，即句jiJ y=O 和 J$\phi$jJy=O( 见图 6.2). 图6.2 等深水域中的表面波 1.等深度水波的基本方程和边界条件的线性化 线性化是解非线性运动方程的有效近似方法，合理的线性化建立在对方程中各 项的量级的正确估计.下面介绍水波问题的线性化方法和步骤. 首先对于水波运动中各物理量作量级估计，假定水波运动在 x ， y ， z 方向的长度 尺度为 ${\rm A}$ ，运动的时间尺度为 T ，而波高 C 的最大值为 A. 假定波高远远小于波长，这 种波称为小振幅波，并有如下的量级估计 z

\[E=?< 1\]

由界面方程 : J\{/Jt+ua${\rm c}$/J .r +viJ \{/Jy-w=O ， 可以估计速度 w 应和 iJ \{/iJ t 同一量 级，即 wOC 乎 在一般水波流场中速度各分量应为同一量级，因而 u-v- 以争 速度势的计算式是 $\phi$= JU. dx ，它的量级应由叫.r的估计推算，即

\[\phi  . c AA -\]

T 用以上的量级估计对各变量作无量纲化 z (.r, y , z ,h) = $\lambda$ ( .r， y ， i ， 的

\begin{gather*}
t = Tt\\
= A
\end{gather*}

$\phi$= 争 (6. 13a) (6.13b) (6. 13c) (6. 13d) 以上各式中带有上标" "的量为无量纲的量，利用以上无量纲公式，我们可有 A). 以及

\[B\phi  T a\phi  Aa\phi \]

iJ.r ). ax T a.i

\begin{gather*}
a 2 \phi  A a\\
2 \phi 
\end{gather*}

iJ.r2 T). a.i2

\[(2\phi A a2 i\theta \]

ay2 n (\}\textasciitilde{}豆 2 (\}2$\phi$A (\}2 岳

\[\theta  Z2 T). ai2\]

将式 (6. 13) 和式 (6.14) 代人水波动力学方程及边界条件，经整理后得 无量纲速度势基本方程 无量纲的水底边界条件

\begin{gather*}
A\phi =0\\
z =-n , - = - = u
\end{gather*}

a" ai 将 z=).z ， \{=A\textasciitilde{} 代人自由面方程 z = \{(.r， y ， t) ， 得 (6. 14a) (6. 14b) (6. 14c) (6.14d)

\[z = e'\]

无量纲的自由面运动学边界条件 豆豆$\downarrow$\textasciitilde{}(\textasciitilde{}全豆豆+坐坐飞=翌 $\theta$ t . -飞 $\theta$ .r J .r ${\rm e}$J y ${\rm e}$J yl ${\rm e}$J z 无量纲的自由面压强条件 苦+?而 -H+EFt=0 第 6 章 *波动力学 线性化是小参量 E$\rightarrow$0 的渐近近似.以上方程中略去 O(e) 的小量，并恢复到有量 纲的表达式，得以下线性水波方程 线化的自由方程 自由面上的运动学和压强条件为 在水底 z=-h 上

\begin{gather*}
A\phi =0\\
z=o
\end{gather*}

\textasciitilde{} \_ Jc${\rm þ}$ Jt ${\rm e}$J z 3$\phi$" 一 +g\textasciitilde{} = 0 Jt O\textasciitilde{}

\[a\phi A\]

Jz (6. 15) (6. 16) (6. 17) (6. 18) (6. 19) 在柯西 - 拉格朗日积分中略去速度势的二阶小量，得 习'古 $\rho$ - P. 一:手+.lfz+ 一一一 o (6.20) dt.. p 式 (6. 15)\textasciitilde{} 式 (6. 20) 是线性化的水波方程和边界条件，方程中的高阶小量都已略去. 特别应注意到水波动力学方程的线性化过程中，自由面 z =e${\rm C}$( .r， y ，t) 的条件已近似 为 z==o. 也就是说等深度水域中速度势的求解域在固定边界 z=o 和 z=-h 之间， 速度势的边界条件和未知量'(.r ,y. t) 无关，这就使求解水波问题大为简化. 2. 等深度水域中小振幅波的解 下面先介绍规则水波的一些基本概念，然后再求解线性水波方程. (1)规则行进波的特征参数 如波面函数 z= ,<.r ， t) 能以三角函数表示:

\[z = acosO = acos( a.r - \omega  t) (6.2 1)\]

则称这种波为单色行避波，其中 a ， Q ， $\omega$ 都是常数 ， a 称为波幅 ， 0= $\alpha$ .r - wt 称为披相 位， $\alpha$ 称为波数， $\omega$ 为披的圆频率或简称频率.波数与波长 A 之间关系为 Q= 草或)，，=主 (6.22)

\[i  Q\]

波数的几何意义为 2$\pi$ 长度上有几个波形，所以它代表波的空间密度;频率和波的周 期 T 之间关系为 向何一 = T 或h-T 一一$\omega$ (6. 23) 频率是 2$\pi$ 时间上出现几个波，所以它代表披在时间上的密度.波的相速度是等相位 点的移动速度，简称为波速，以 r 表示.由式 (6.2 1)给出的单色行进波具有线性相位 函数。 =$\alpha$.r -wl ， 令 d8=$\alpha$dx 一 $\omega$dl=O ，可得波速 z = (生) . $\tilde{ft}$.. -旦 =$\lambda$dl 'O \textasciitilde{} 常敬 $\alpha$ T (6. 24) 如果披速(披的相速度)和波长或频率无关，这种波称作非色散波:如果披速和波长或 频率有关，则称为色散波.下面将会看到深水波属于色散波. 波高的极大值称为波峰;波高的极小值称为波谷.波高函数等于零的点称为波 节点. 任意个单色行进波的线性叠加称为多色波，故一般的规则多色波可表示为 z- 衍，L) - L:a;cosO, L: a;co 巾;x 一 $\omega$川 (6.25) 如果 ${\rm C}$(x ， t) 是平方可积的任意函数，则利用傅里叶积分将它分解为连续分布的 多色波: z = ${\rm C}$(x ,t) = 去j 以川 (6. 26) a, ( $\alpha$) 称为波谱.对于线性的小振幅水波来说，它的控制方程是线性的，只需要研究单 色规则波的传播特性，多色波的传播可以用单色披的线性叠加来分析. (2) 等深度水域中规则行进波的速度势解 写出水波的速度势方程和边界条件.

\[A\phi =0\]

z = 0: 窍 -2=o z = 0: 旁 +g${\rm c}$ = 0 . a $\phi$ 《 z =-n: 王三 -u 假定水波具有 ${\rm c}$=acos($\omega$.-wO 形式的行进波解，则由式 (6.27b) 得 (翌). \textasciitilde{} " =$\omega$in( a.r - wl) (J z ，. \textasciitilde{} 。 因而可以假定在 -h"\$:二 z<O 的水域中

\[\phi = c\theta ( z) sin( a.r - wl )\]

代入式 (6.27a) 得 (6.27a) (6.27b) (6.27(') (6.27d)

即 它的解为 (:2 ) i$\tau$ -a2$\tilde{lsi}$n(ar -wt) = 0 oz- I d2\textasciitilde{} \_ ..2 ..1. dz 2

\begin{gather*}
\tilde{r} -\\
~ = CI expaz + C2eXp(- az)
\end{gather*}

第 6 章水波动力学 由水底边界条件式 (6. 27d): (a\textasciitilde{}/aZ). =\_ 1I =0 ， 因而 CI exp( - ah) - C2 exp($\alpha$h)=O ， 于 是有: CI =c2exp2叫，消去 CI 后，\textasciitilde{}的解可整理为 \textasciitilde{}= :\_2 expah expaz - +czEXp( 一回)

\[exp(- ah)\]

\_.. exp$\alpha$ (h 十 z) + exp[ -a(h+z)] q exp( ah)

\[=-JEL-chU+z) exp(- ah)\]

于是水波的速度势等于: $\phi$= Cch$\alpha$ (z + h )sin(ar 一 $\omega$ l ) 式中 C= 2C2 exp(ah). 将速度势代人自由面条件(式 (6.27b)) ，得 :C$\alpha$ sh(ah) =a $\omega$ 或 C=a$\omega$/[$\alpha$ sh(ah 汀，从而行进波 !;=acos($\alpha$.1' -w/) 的速度势为 $\phi$= --'坦一7cha (z + h)sin($\alpha$X -w/)

\[\alpha sh(ah )\]

(3) 等深度水波的色散性 (6. 28) 水波运动必须满足自由面处压强边界条件， $\theta$$\phi$/at+g!;=o ， 把 $\phi$ 和 C 代入该式 可导出水波运动的波数和频率之间的重要关系式 z 整理后得 二坐三 ch(ah) + ga = 0 ash(ah ) 旷= ag ${\rm •}$ th( 础) (6.29) 式 (6.29) 是表面波的重要性质 z 称为色散关系式.由波速公式 C= $\omega$/$\alpha$ ，可导出线性水$\cdot $ 波的披速公式 z C=\textasciitilde{}乎 (6.30) 可以看到线性水波的波速和波数有关.波速和波数有关的性质称为色散性，对于单色 波，色散性并不表现出来.多色波，即由若干个单色波叠加的组合波，每个单色波在行 进过程中的波幅、波数、频率保持不变.如果这种波是非色散的，即不同波数的单色波 以同一波速行进，因此，各个波之间的相位关系不变，初始的披形在行进过程保持不 变.如果波是色散的，不同波数的单色波行进速度不同.各个波之间的相位关系发生

变化，组合波的波形在行进过程中不断地变形.稍后将详细讨论水波的色散性. 如果水深很浅 .h $\rightarrow$ O. 这时 th($\omega$ $\rightarrow$川，因此波速 c= ，Jgh. 它和波长无关，就是 说等深度浅水线性波是非色散的.如果水深很深 ${\rm •}$ h $\rightarrow$$\infty$，这时 [th (ah ) Jh…= [sh(ah) / ch(ah) Jh-" = 1. 因此波速 c=\textasciitilde{}=J百万言，这时波速仍然和波数或波 长有关.就是说无限深水中线性水波也是色散波.将 g=9.81m/s 2 代人披速公式，得 无限深水中波速的简单公式 z

\[c = 1. 25 ..5\]

所以，无限深水中波速和波长的平方根成正比，也就是说，无限深水中长波的相速度 大于短波的相速度. 下面先来考察水波运动的流场性质. (4) 等深度小振幅水波流场的流线和质点轨迹 由速度势 $\phi$ 可求得速度场为 u= 坐=一丝$\tilde{cha}$(z + h)cos(a .I - wt)

\begin{gather*}
<JI sh(ah)\\
J\phi 
\end{gather*}

w= 一=一旦旦;-;-sha(z+ h)sin( a.r -wt)

\[<J z sh(ah)\]

${\rm 1}$流线 给定瞬时的流线方程为 dx/u=dz/切，即

\begin{gather*}
dz sha(z+h)sin( a.r -wt)\\
dx cha(z + h )cos( a.r - wt)
\end{gather*}

积分后得 (6.31a) (6.31b)

\begin{equation*}sha(z+h)cos(ax -wt) = const.\tag{6.32}\end{equation*}

由式 (6.31a) 和式 (6.31b) 可知在单色行进波的波峰和波谷处 (COS(ax 一 $\omega$ t) = 士 1 ， sin($\alpha$ x-wl) =0). 水平速度数值最大，垂向速度 w=O; 而在节点处 (cOs(aI-wt) =0. sin(ax-wt) = 士1).水平速度等于零，垂直速度数值最大. ${\rm 2}$水波中的质点运动轨迹 轨迹方程为 生 =u= 一旦$\tilde{cha}$(z+ 的 cos( a.r一 $\omega$ t)dt - sh( 础) 生 =w= 一旦:-:-sh$\alpha$ (z+h)sin(aI -wt) dt - sh($\alpha$h) \_.\_-- ,- , ... ..\_--,\_ 质点轨迹的做分方程的精确积分不能用初等函数表示.对于小振幅水波，流体质点在 原静止平衡位置附近振荡，因而令方程式右边 z=zo .x=xo 为常数，求眩轨迹的近似 积分.得 骂王 = a ， wcos( $\alpha$I o 一 $\omega$仆， dt 生 = h ， $\omega$sin( a.r o - wt) dt "'W"---, \_U

式中: al =a ${\rm •}$ cha(ZO +h)/sh(a的 ， b 1 =a ${\rm •}$ sha(ZO +h)/s h( a的，最后得 x - xo = alsin[旺。一 $\omega$ (t - to 汀 ， z - Zo =- b1 COS[axIY 一 $\omega$ (t - tO)J 消去 t 后得

\[(X-X O) 2 I (Z-Z O)2 - 1\]

\textasciitilde{}'一亏了一-.1. 即水波中流体质点的近似轨迹为椭圆

\begin{gather*}
a. = ~. cha(z o + h ~>~. shaz o + h~=h\\
1 - sh(ah) ./ sh(ah) - <11
\end{gather*}

(6.33) 故 al 为椭圆长轴 ， h 1 为椭圆短轴.就是说单色线性水波中流体质点的近似轨迹是长轴 在水平方向的椭圆.在水底 Zo =-h ， 得 b 1 =O ， 这时椭圆轨迹退化为直线，这是由于 水底壁面条件的限制，流体质点只能平行于壁面作直线振荡，如图 6.3(a) 所示. ' 一

\['/////////// / /////////////////////\]

(a) 另立 I (b) 图6.3 线性行进水波的质点运动轨迹 (a) 线性行进水波流场中的质点轨迹示意图; (b) 线性行进水波流场中的质点轨迹流动显示图 图 6.3(b) 是线性行进水波流场中质点轨迹的流动显示图像，在水波流场中撒入 细微颗粒，颗粒的质量密度和水很接近.曝光时间为一个披周期，就可以摄下水质点 的轨迹，如图 6.3(b) 所示.图中清楚地显示质点轨迹的椭圆度从水面以下逐渐增加. (5) 无限深水中的行进波 当水深 h 为无限大时，表面行进波的解可以由有限深水波解求极限得到，也可以由 直接求解水波方程得到，为简单起见，用有限深水公式来导出无限深水中水波公式. 「 $\omega$ch[a(z+ 的 J\_:\_ ， c \_ .\textbackslash\{\} l $\phi$ (x ， z.t. $\infty$) =liml \textasciitilde{} \textasciitilde{}"LU.'\textasciitilde{} .' ,"'..J sin(ax -oot) I h阳 L a sh($\alpha$h) \textasciitilde{}...，- \textasciitilde{}-， J 一 $\omega$1:\_r expa(z+ 的 +exp(-a(z+h))lI-- YW '- 'J..:"" ""....Y'" ""'\textasciitilde{}，，' .." Isin(ax- CtJt) ah:: L exp(ah) - exp( - ah) r...\_..... W =型 exp($\alpha$:z)sin(ax -oot) a

即无限深水中的速度势为 此外 在无限深水中色散关系为 $\phi$= 但 exp(az)sin(ar -wt) chah limthah =一一一= 1 h阳....\_- shah (6. 34)

\begin{equation*}\omega Z=ga\tag{6.35}\end{equation*}

元限深水中水波的速度场为 u= $\omega$exp($\alpha$z )sin(ar 一 $\omega$ t). w = $\omega$exp($\alpha$z )cos( a.r - wt) 速度场流线可用上面类似方法求出: exp($\alpha$z)cOS(ax 一 $\omega$ t)=const. (6.36) 流体质点的轨迹仍用近似方法求它们在平衡点附近的运动轨迹，有 生 =$\omega$ exp(azo )sin(axo - wt) 生 =$\omega$exp($\alpha$'zo )sin(aro - wt) dl - UU/":;\textasciicircum{}tJ" U'""'O ，f $\tilde{lU}$' U'.A. 0 WL" , dt - UUJ \textbackslash\{\}:,:. \textasciicircum{}P 积分后得

\begin{equation*}(X-X O) Z +(Z-Z U)2 = (aexpaz o)2\tag{6.37}\end{equation*}

就是说无限深水波中流体质点轨迹为圃，回半径随水深以指数函数减小. 从以上讨论可以看到无限深水域中的小振幅水波的运动和流场性质与有限深水 波本质上是相同的，但无限深水波公式比较简单，如果有限深水波能用元限深水近 似，则运算要方便得多.通常以色散关系作为近似的依据，这两种情况色散关系相差 的因子为

\[\omega 2/ ga = th(\alpha /1)\]

由双曲正切的公式可得 th(2$\pi$)= 主芋$\zeta$= 0.9999932 句 1 e-" 寸- e -. 也就是说.当 h> 主=).时，无限深水的近似在色散关系计算中的精度达 10 - 6 ，因此 通常水深大于被长情况下的等深度小振幅水波 可以用无限深水波近似. (6) 等深度水域中的小振幅驻波 上面讨论的行进水波，在行进方向，即 z 方 向，没有边界限制，表面波可以自由地在 z 的 正、负方向传播.如果研究矩形容器中的水波时， 在容器的垂直壁面上，流体的 r 方向速度必须 等于零(参见图 6. 的，这时，应附加边界条件: 图6.4 矩形容器中的驻波

218 第 6 司E 水波动力学

\[L ~-L eJ \phi \]

当 x= ::t一时.一 =0\textasciitilde{} 2 \textasciitilde{}.. ${\rm e}$J x 于是水波的基本方程和边界条件为

\begin{gather*}
A\phi =0 (6. 38a)\\
3\phi \theta \phi \\
z = - h: .~ ~ = 0 , x =:::!:: L: .~ ~ = 0
\end{gather*}

。 z ox (6. 38b) = 0: \textasciitilde{}全-\textasciitilde{} = 0 , \textasciitilde{}壁 +g${\rm c}$ = 0 ${\rm e}$J z ${\rm e}$J t - , Jt (6.38c) 由于速度势基本方程和时间无关，它有分离变量解 $\phi$= 骨 (x ， z)sin$\omega$t 代入基本方程，得 的边界条件为 ZJ. =-h ， 一= O. -Jz'

\[nu --\]

忡 -hL $\pm $ 一Z 用经典的分离变量法解拉普拉斯方程，令 骨 = X(x)Z(z) 有 d2 X d 2 Z( 到， 一-.-.- =一一寸一=矿 >0dxZ dz${\rm ‘}$ 以上方程的解为

\begin{equation*}X(x) = C1 sinax + C 2 cosax. Z(z) = DI expaz + Dzexp(- az)\tag{6.39}\end{equation*}

水底边界条件 (当 )=0 ，即 (dZ/ 由)=\textasciitilde{} 10 = 0 ${\rm e}$J z J. $\tilde{h}$ 容器侧壁边界条件 (坐 )=0 ，即 (dX/ 创$\iota$=0${\rm e}$Jx J J \textasciitilde{} :!:L 将水底边界条件代入式 (6.39) ，得 D1exp($\alpha$ h) - D2 exp ( - ah) = 0 ，消去一个任意常 数后 ， Z(z) 可简化为

\[Z(Z) = Dch[a(z+h)J\]

函数 X(x) 中的常数可由侧壁边界条件求出 C1cos 等- Czsin 等= 0 , C1cos 等 +C2 5 咛=。 进行简单的代数运算可导出 C1 = 0 和 aL n

\[5m 2 = u\]

即 $\alpha$L \_ 2n $\pi$ 」 $\tau$ =n$\pi$ 皿 $\alpha$=$\tau$- 现 ${\rm A}$= 主 ， n = 1, 2,… n (6.40) 这就是说，有界容器中驻波的波长由容器的宽度决定，不同的 n 值是谐波的等级. 将 X(.r) 和 Z\{Z) 代入速度势解公式，合并任意常数后，得 $\phi$ = Ccosaxcoswt ${\rm •}$ cha(z + 的 由水面边界条件J$\phi$/Jt+g${\rm c}$=o ， 得水面方程为

\[= cos\alpha \tilde{xsl}n\omega t .-g\]

令 C$\omega$ch\{ah)/ g=a 为驻波振幅，则驻波解为 $\phi$= \_\_\_\_E\_丘一7cha( z + h )cosaxcos缸， t

\[\omega ch(\alpha h)\]

(6.41a)

\begin{equation*}c = acos\alpha 'XSln\omega  t\tag{6.41b}\end{equation*}

驻波的波形如图 6.4 所示，单色驻波的波形始终保持余弦形..:r= :1:1- /4 始终为节点， .1 '=0 和 x= :1: L/2 始终是极值点;而驻披的幅值随时间变化 asinwt. 最后由自由面运动学条件J$\phi$/ ，)z -a${\rm c}$/rJ t=o. 得色散关系: $\omega$2=$\alpha$'g ${\rm •}$ th\{$\alpha$I!) 它和行进波色散关系相同. 驻波的速度场为 --二旦旦\textasciitilde{} , cha(z + h )sinaxcoswt. 'W 一一旦旦旦--:-sha(z + h ) cosax coswt $\omega$ch(ah)\textbackslash\{\}"'llU'" I H .I.;J I J.l U..(. ，"" V o71 W I..一 $\omega$ch'( aIz) s 由速度表达式可以看出，驻披的波峰和波谷处( COSax = :1: 1 )水平速度等于零，垂直速 度的数值最大;而在节点处 (COsa .r=士们，垂向速度等于零，水平速度数值最大.将速 度场代入流线方程

\[d.r - dz\]

cha(z 十 h) sina.r sha( z + h) cosax 职分后得流线谱为

\begin{equation*}sha(z + Iz )sinax = const.\tag{6.42}\end{equation*}

和行进波的流线谱做对比(式 (6.32): sh$\alpha$ (z+ 1z )cos\{$\alpha$ .r -wl) =const. )，给定瞬间单 色驻波的流线谱和单色行进波的流线谱相同，相对于波面有 $\pi$/2 的相位差.因为，对 于驻波来说，波峰和波谷处的质点垂直速度总是最大(绝对值) ，因此流线垂直于波 峰;而行进波波峰处的质点垂直速度始终等于零，因此行进波中的流线谱和波峰及波 谷相切.驻波的流线如图 6. 5 所示，它是用流动显示方法得到的图像，和式 (6.42) 完 全一致.图中每一线段代表质点的一段轨迹，和图 6.6 的计算结果一致. 质点在平衡点附近的近似轨迹为 生=二坐巳 cha(z" + Iz ) sinax" coswt

\[dt \omega ch(\alpha Iz ) \upsilon \upsilon \]

积分得 生=一旦$\tilde{sha}$ (z" + h ) cosa.r n C()SWI dl$\omega$ch( aIz ) ..\_-\_.-, 第 6 章水波动力学

\begin{equation*}x-xl) =alsin\omega  (I - 111 ). z - z" = b l sin\omega  (1-/ 0 )\tag{6.43}\end{equation*}

式中 al=-二旦旦$\iota$ cha\{z" +h)sinax" .b, = $\tilde{I}$!a.. sha\{zu +的 coSQXo. 闲此.驻波中流$\omega$ ch($\tilde{h}$) cna$\tilde{z}$" Tn JSlnaX" .IJI = wc盯否? 体质点在平衡点附近作直线振荡. ${\rm ‘}$ ha\{zo 十 h)co姐ro 仨去=去=叫，叫=山 (zo + h)sin a.ro 质点振荡的振幅和斜率随水深和驻波的相位变化，如在水底 z= 一儿，质点轨迹的斜 率叫=0. 质点作水平振动，其振幅为|品和叫在波峰和波谷处 :X=O 及 X= 土 L/2. 质点作垂直振动，因此 tan${\rm ß}$= 士$\infty$ z 在节点处 : aX= 士 $\pi$/2 或 X= 土 L/4 处，质点作水平振动，因此 tan${\rm ß}$=O. 驻波水面附近的质点轨迹如图 6. 6 所示，它和 图 6.5 显示的水面附近以直线震荡的结果$\rightarrow$致，因此线性水波的近似是合理的. 图 6.6 驻波水面附近的质点轨迹 实线 g 波面 s 虚线 z 质点轨迹 (7) 行进水波流场中压强 由线性近似的柯西-拉格朗日积分式得 $\rho$ 一户. \_ il $\phi$ P ill 口 将速度势解 (6.28) 代人上式，并利用色散关系式旷 =ag ${\rm •}$ th(a的，便得行进波中水 的压强 z 主二生=一旦旦;-;-cha(z + 的 cos( a.r - wt) - gz p ch\{ah) (6.44a) 或

\section*{6.3  线性水波的色散关系}\addcontentsline{toc}{section}{6.3 线性水波的色散关系}

主二主=一互$\zeta$"7 cha(z+h)- J!z (6.44b) p ch($\alpha$h ) \_..\textasciitilde{} ,- , .., ... 可以看到，水波流场中的压强由两部分组成，一部分是静水压强(- gz). 它和时间无 关 z 另一部分和水波表面高度有关，这部分压强也是以行进波方式在水中传播.

色散关系是重力场中表面波的重要性质，本节着重研究水波的色散性.以无限深 水域中水波中的色散关系为例: 也\textasciitilde{}. = ga 波速 c=:=\textasciitilde{}=${\rm Æ}$ 对于单色波 ${\rm c}$=acos(a.1 一 $\omega$t) ， 它的波形在行进中始终保持余弦曲线，只是余弦函数 的相位以等速度 c= $\omega$/$\alpha$ 移动.当考察多色波时，由于长波相速度大于短波相速度，波 形(波面函数的几何形状)因色散而发生变形.现在考虑一种最简单的多色波 z 两个 振幅相同、波数不同 ($\alpha$l 手 $\alpha$2 )的单色波叠加 z ${\rm c}$ = acos($\alpha$l .r一 $\omega$ It)+acos( 的 I 一剧。 这时两个单色波的相速度是不同的: 1= 去 =${\rm Æ}$=J言， tz=TJZ= 在 由于相速度之差，两单色波的相位差随时间不断变化，从而波形发生变化.由于两单 色波的相位差不断改变，直觉观察到的被高最大值的传播速度不再等于单色波的相 速度. 1.波包的群速度 两个振幅相同、波长十分接近的单色行进波叠加称为波包.波包的波面方程为 ${\rm c}$ = acos( a.r - wl ) + acos[ ($\alpha$ + aa ).r 一 ($\omega$+a$\omega$ )tJ 现在来考察波包的波形演化.由主角公式: cosA+cosB=2cos[(A + B> /2Jcos[(A

\[B> /2].\]

当 8$\alpha$ <<a 时， 或 ${\rm c}$ = 2ac 叫 tsar- 护$\omega$ I ) 叫 ($\alpha$+ 护).r一 ($\omega$+ 学 )tJ C=2um48$\alpha$l .r一坠t)cos( a.r -wt) $\iota$ 飞 oa ,

\textasciitilde{}=巾 ， t)c以u 一 $\omega$ 仆 ， a( .r, t) = 2ac叫护(.r一安 )J 上式的运动学意义为 z 波包的振幅不是常数，它本身以 8$\omega$/ '8$\alpha$ 速度行进.所以波包属 于波幅调制的行进波. 定义 6. 1 色散波的频率对波数的导数称为波的群速度，或披群速，并用 ck 表示 z

\[d\omega \]

Eu= 百二 (6.45) 对于波包来说，群速度是波振幅轮廓的传播速度.图 6.7 中实线是波面，虚线是波面 极值的包络，包络以群速度 C. 运动. 《miZM立YX》 图 6.7 波包示意图 实线 g 波回 B 虚线=波包

\subsection*{例 6.1  求有限等深度水域和无限深水域中线性重力波的群速度.}

解有限深水的色散关系为

\[\omega 2 =guth(\alpha >h)\]

求导数 d$\omega$/ 白，得求有限等深度水波的群速度为 c. = 坐=丘 th(ah) + 一骂$\iota$一da 2$\omega$Z$\omega$ch品 (ah ) 令 h $\rightarrow$$\infty$ ，得无限深水波的群速度为 c.. = 坐=丘= \_!\_ /互 =ic $\tilde{d}$$\alpha$2$\omega$2 勾 $\alpha$2 无限深水中线性重力波的群速度为相速度之半. 2. 波能及波能的输运方程 定义 6.2 单色波单位长度内的流体平均动能和势能之和定义为波能. 首先计算一个波长的波动能，应用第 5 章不可压缩无旋流动的动能公式 z T=ffI$\phi$ 芳出 式中积分周线为:*面 z=O ， 水底 z=-h ， 左边界.r =0 ， 右边界.r =${\rm A}$=2$\pi$/$\alpha$ .n 指向积分域外(见图 6.8 )，根据法 向量规定，在 OC 面上 J$\phi$/Jn = - ${\rm e}$J $\phi$/CJ 趴在 AB 面上 3$\phi$/

\[\lambda x\]

Jn= ${\rm e}$J $\phi$/ ${\rm e}$J.r;同理，在 OA 面上J$\phi$/Jn= $\theta$$\phi$/ ${\rm e}$J z. 在 BC 面上 C B a$\phi$/$\theta$ n=- ${\rm e}$) $\phi$/ ${\rm e}$) z. 于是动能积分式可写作 图 6.8 推导波能月1 阁

T =?引扎t九$\mu$t州114 ($\phi$ 芳)九宫FJ=40yd由5 +jL汕 (归$\phi$ 芳 L =卢叫JAd由s+L肌 (伊$\phi$ 尝)九::-=一=$\leftarrow$二一「叩${\rm ‘}$$\infty$卢? =f刊[f川I:(仰$\phi$ 尝).=0 dx+ t \textasciitilde{}$\infty$ (但$\phi$ 主)$\lambda$Z户叫=吐A 缸 +I: 一 (归$\phi$ 尝凯).$\lambda$zo = .c .... 产世 +L'二飞一寸(但$\phi$ 兰t九\$=0 dz ] 在无穷远水底 ($\omega$ $\phi$/川$\theta$ z) 、$\upsilon$ =0 ，困而第三项积分等于零;由于单色波的周期性， ($\omega$ $\theta$ $\phi$/川rJ.r川) .，=斗0 = (rJ $\phi$/J .r):r = ;. ， 因此第二项与第四项积分和等于零.于是波动能积分式简 化为 T = $\tilde{I}$: ( $\phi$ 主) .=0 d.r 无限深水域中线性行进波的速度势 z $\phi$= 纽 exp($\alpha$z)sin( a.r -wt) 由此得 $\phi$I. =u =纽 sin(a .r -wt\} ， ( 坐) =旦旦 sin(a .r一 $\omega$t) ， 代人动能公式 z$\omega$ '\textbackslash\{\} rJ zJ.= o $\omega$ T= 旦豆华 r;. sin2 (a.r - wl ) d.r

\[2\omega  ' Jo\]

注意到无限深水的色散关系 $\omega$2= 附和I: sin2 (a.r一川.r =$\lambda$/2 ，积分结果得 T= 丘 a 2 g). 4 .. (6. 46a) 再计算一个波长内的表面波势能.水面高为仁宽为 d.r的水柱的质量为 pg$\tilde{d}$.r， 重心位置为 \textasciitilde{}/2 ，故势能为 pg$\tilde{d}$.r /2 ，一个波长内的势能为 E=j:$\div$阳 \textasciitilde{}2 将 \textasciitilde{}=acos(k .r -w/) 代人上式，得 II= 拉a 2 ). (6. 46b) 在无限深水域中线性行进披的波能(单位氏度内平均动能和位能之和)为 a 阳1-2 一一 E 一 +一、 A

\begin{equation*}T--E\tag{6.47}\end{equation*}

以上计算表明: (1)线性水波的波能中动能和重力势能各占二分之一; (2) 披能和波 幅的平方成正比. 下面考察水波行进过程中能量的传递.在行进波中任取一垂直截面，则在一个时 间周期内该截面上压强所做功应等于向截面另一侧水波输送的能量，截面左侧流体 压强对右边流体所做功的功率等于 z

\[w'= r\rho \]

将行进披中压强公式 (6.44) 和流向速度公式 (6.31a) 代入并积分后，得 单位时间内压强做功等于 w'= 旦望

\[2\alpha \]

w= 些， $\zeta$庄E T 4 式中 c= .;;;J再是波速，将波能 E=pga 2 /2 和群速度 c. =c/2 代入上式后，就有

\[w = \& '. (6. 48)\]

上式是单色波巾被能的传输公式.它表明，无限深水域的单色波中，压强单位时间向 行进波传递的能量等于被能和群速度的积.简言之，无限深水域的单色波中披能以群 速度传播.上述结论对于等深度单色波也成立，读者可将等深度水波的公式代人上面 的推演公式得到证明. 3. 波阻 物体在有自由面的水中运动时不断激励向外传播的表面披，因而物体必须不断 向水波提供能量.根据功能原理，物体必须对流体做功，同时流体对物体有反作用力， 这种流体对物体的作用力称为兴波阻力，简称披阻.利用披能公式和波能传输公 式 (6.48) ，可以估算简单运动的披阻.如有二维物体在有自由面的无限深水中运动速 度为 u ， 因物体运动而兴波，设 F 为克服兴披阻力作用在流体上的合力，则物体对流 体做功为 Fu. 船只在深水中等速 u 直线稳定行进时，在船体前形成的水波以波速 c=u 向前传 播.设该水波的波能量为 E ，则水波不断向外传播能量 uE. 但该能量中有 一 部分是波 场中通过压强作用传递过来，由式 (6. 48) 可算出这部分功为 c.E ， 故需从外界输入的 能量为 (u-c.)E ， 根据功能原理，应有 "'*- r::- \_ (u - c. ) E Fu = (u - c., )E 亚 F= 一一一一一一 u 在无限深水中 C. =c/2 = u/2 ， 代人上式，得二维行进波流场中，物体所受波阻

\[F=5\]

如果已知物体激励的水波波高 a ， 曲式 (6. 47) 可以计算波能 E=pga 2 / 2 ， 从而算 出波阻. 一 般 三 维物体在水中运动时受到的波浪力应通过求解水波方程求出速度势， 然后应用波能公式，或计算物体表面的压强来计算波阻.有关内容请学习水波动力学 的专门著作或文献.

\noindent\textbf{6.4.} 缓变水深申线性水波的传播 22:;

\section*{6.4  缓变水深中线性水波的传播}\addcontentsline{toc}{section}{6.4 缓变水深中线性水波的传播}

上面讨论了等深度水域中的波传播，实际中常常有变深度的水域，例如，海岸或 浅滩都是变深度的水域.由波速公式 c = $\omega$ /a= ..Jgth(ah)/$\alpha$ 可以推断，当水波由深水 区向浅水区行进时，传播速度将减小.依照波动现象的 一 般原理，当波从 一 个区域进 入另 一 个 K 域发生波速改变时，将产生折射现象.本节主要从能量传输角度研究变深 度水域中的水波传播性质. 下面讨论变深度水域中的波传播.为了讨论简单起见，假定水深是缓变的，即水 底方程可写作: h = h( $\mu$x) ${\rm •}$ $\mu$<<1 ，更具体地，引人如下的小参量: '1 = o( 但生)

\[\alpha h\]

(6.49) 式巾分子是水底斜率，分母叫可以表示水波的波面陡度.它和水面的最大斜率同一 量级.式 (6.49) 定义的参最严表示水底斜率远远小于波面陡度.下面用摄动法导出水 披的折射方程. 1.缓变水域中的波能传播方程 可以想象，变深度水域中的线性水波的波数和频率将发生变化，就是说水波在变 深度水域小行进时，它的波数或当地的水波频率将发生变化.当水深是缓变时，波特 性也是缓变的，即 :$\omega$=$\omega$($\mu$ x. $\mu$ 1) .$\alpha$ = $\alpha$($\mu$h$\mu$ t) 等.缓变的水波波数和频率特性的含 义可直观地解释如下:观察者在 一 个或几个波氏上看到的是均匀的水波，要在很长 的距离上(比水波的波长大 一 个挝级)才能观察到水波的波长和频率发生变化.换句 话说，在实际的物理长度和l 时间上，可以定义局部的服数和频率.通常将 x= $\mu$x.t = $\mu$t 称为缓变量或慢变量 ; .1" .1 为快变量.在 一 个或几个波长上水波的色散关系不变， 仅和当地的水深有关.即局部的色散关系不变:

\[\omega 2 = agth(ah)\]

频率和波数部是慢变量的踊数 :$\omega$ = $\omega$ (x. 仆， $\alpha$=$\alpha$( 土 .1 ).如果以 dx/dt = $\omega$/$\alpha$ = c( .i ${\rm •}$ 1 )表示水波的轨迹.这时波的相速度不同是常数，而是流向坐标和时间的缓变函数， 也就是说波射线不再是直线.在缓变水域中.我们定义当地的波相位角为 ()=$\alpha$ x wl. 则略去高阶小量后.有 d() = $\alpha$ (.1'. t ) d . r 一 $\omega$ ( .1' .Odt + ()($\mu$) = ($\alpha$d i" 一 $\omega$d l) / p = d()/$\mu$(6.50) 式 'I '${\rm O}$(x.t) 是慢变的相位面数.应用求导数公式，应有 z$\alpha$ =JO/Jx=J${\rm O}$/(]X 和 $\omega$= - r'J()/ rJt = - a${\rm O}$/rJ l ${\rm •}$ 以及 a 2 () / iJ x r1 1 = $\theta$ \textasciitilde{} ${\rm o}$ / al iJ.1.'. 于是很容易导出缓变水域中波特性演 化方程:

\[3\alpha  ( .r,  n , a\omega  ( .r,  t) \wedge \]

干-.. r1 r J .r (6.5 1) 式 (6.51) 可以理解为波相位的连续方程.在缓变坐标系中 $\alpha$ 是"波密度"(单位长度中 波的个数) ， $\omega$ 是"波流量"(在固定点上，单位时间流过波的个数) .如果用 $\omega$=$\alpha$c 代人 式 (6. 5 1)，则它和流体力学中的连续方程十分相像.由于波的色散关系旷= $\alpha$gth(a的，式 (6. 51) 还可进一步简化为 ${\rm e}$Ja( 3\_' ,f) I d$\omega$ 句 (.i，[)\_iJ$\alpha$ ( .r， 。${\rm ‘}$ $\theta$$\alpha$ ( .r， 1) \textasciicircum{} 二千一一一一千 f 一一 J1 'da i1.r 一缸 '-Hat w 由式 (6.52 剖，还可以导出=

\[3\omega  ( .r,  n I iJ \omega  ( .i, [) \wedge \]

-----'---r- c.. = u Jl ${\rm •}$ J:i (6. 52a) (6.52b) 式 (6.52a) 和式 (6.52b) 表明，沿着群速度轨迹缓变水波的波数和频率不变.这一结果 说明缓变水域中波传播的运动学特点.在等深度水域中单色波行进时，在相速度轨迹 上总是能观察到该单色波的披长和频率.但是，当单色波进入变深度水域时，由于水 波特性的变化，在当地相速度的轨迹上水波的波长和频率是缓变的，只有当观察者以 群速度运动时，他才能观察到同一波数和频率的波运动. 2. 缓变水域中的水波速度势的摄动展开 下面来导出缓变水域中波传播的动力学关系式，基本思想是用缓变量做摄动展 开.该方法来自量子力学中著名的 WKB 方法 (Wentzel-Kramer-Brillion) . 由于水底是缓变的，可写成 h=h(p抖，这表示水深是波传播方向上坐标的缓变 函数.而在水深方向，水波流动的特征长度并没有变化，因此，引人新的坐标量

\[.r =\mu r ,  z = z , t = pt\]

为了分析的方便起见，设定复数形式速度势: $\phi$=$\phi$ ${\rm •}$ (:i， 坛 ， t ， $\mu$)expi8/$\mu$+c. c. (6. 53) ${\rm O}$=\{J(:i , i ，[)是缓变的相位函数 .$\phi$$\cdot $ 是缓变的幅值函数， C. C. 表示复共辄.幅值函数 旷是含小参量 $\mu$ 的函数，它可进一步用小参量 $\mu$ 展开为 $\phi$ . (:i， 主，[) = 1>0 + (一 i$\mu$ 冲 1 + (一 i$\mu$) 咱2+ … (6. 54) 以上摄动展开的思路是 z ${\rm 1}$7.1<波幅值是缓变的;${\rm 2}$水波特征 $\alpha$ ， $\omega$ ， C 等也是缓变 的.下面导出变深度水域中幅值公式. 3. 幅值演化方程 将摄动展开式代人线性水波基本方程和边界条件:

\begin{gather*}
A\phi  =0 ,  -h<z<O\\
3\phi 3\phi  ah\\
z =-h( .i): .:~+.:~ :" = 0
\end{gather*}

。 z rJ.r rJ.r

z = 0: 尝 +g${\rm c}$ = 0 和 去一窍 =0 最后一式也可写作 z

\[i:J Z \phi 8\phi \]

z = 0: \textasciitilde{}....;- + g 一 =0${\rm i}$:J t' . \textasciitilde{} ${\rm i}$:J z 用式 (6.50) 可将 ${\rm i}$:J/${\rm e}$J:r，${\rm e}$J/ $\theta$t 等变换为

\[3-M\]

$\mu$ 一一 -r-r 吨。-吨。主任 一

\[a-h eJ i:J ai i:J\]

at 否 37=$\mu$ ${\rm e}$J I' ${\rm i}$:J ${\rm i}$:J

\begin{gather*}
32-32'\\
32-282
\end{gather*}

石言一 $\mu$ 否? 将它们代入基本方程和边界条件中各项，例如 3$\phi$$\theta$$\phi$3$\phi$. , . \textasciicircum{} I '- I ${\rm •}$ ao 了 =$\mu$-=$\mu$-T-exp(iO/$\mu$)+It 旷 exp( i\{\} /$\mu$) 曰 X r ${\rm u}$:r r' ${\rm i}$:J${\rm I}$ . -. r- iJ${\rm I}$ a Z $\phi$ .， ${\rm i}$:J z $\phi$ ? a z $\phi$ ${\rm •}$ . . \textasciitilde{} . . \_. a $\phi$ ${\rm •}$ a1J , .:;; 石2" -$\mu$" 苟言 =$\mu$Z 7ZTEXP(i0年) +2i$\mu$7F 再 exp(i8/$\mu$)

\[iJ21J ~ . , .;, I ,  1 \theta 1J   Z\]

+i$\mu$ 石E 旷 exp(i1J /$\mu$) (石 )$\phi$ . exp(i8/$\mu$) C] 2 $\phi$ ? rJ 2 $\phi$ ${\rm ‘}$ :;. ， , . \_ . a $\phi$ ${\rm •}$ ${\rm e}$J1J -:;-:2 =旷一.$\tilde{Z}$ exp(i8 /$\mu$) +2i$\mu$-一 :V exp (i8/$\mu$)rJ t' ,- rJ t2 ---\textasciitilde{} ' -- ' '-' . --,- ${\rm e}$J t at aZ1J \textasciitilde{} ${\rm •}$ , .;, 1 , ${\rm e}$J 21J +i$\mu$ 言\_ ， Z $\phi$. exp( i8/$\mu$) 一一$\tau$$\phi$. exp(i8/$\mu$) 二 rJ l

\[\theta \phi 3\phi \cdot \]

一二一$\tilde{exp}$(i8/$\mu$) , ${\rm i}$:J z az

\[eJ 2 \phi \theta 2\phi .\]

一一 =-77exp(i0年〉iJ z Z ${\rm i}$:J z 将变量置换后的公式代人基本方程，并按(一 i$\mu$)0 ， (一 i$\mu$) I , (一 i$\mu$)2 同量级项归并，得 以下方程组. (1)零阶近似.量级 O( 一 i$\mu$ 沪的基本方程及边界条件如下 z 骨';-az 骨。 =0 ， . -h< 运 <0 nu ，由 rnu 2

\[\omega -K\]

一一

\[,PE --\]

一一

\begin{gather*}
-Z\\
-Z
\end{gather*}

上标撇号表示对 z 求导，不难看出，零阶近似就是等深度线性水波，其速度势解为 (困 z= 玉 ，下文取消上标横杠) : 仇=\_ igA $\tilde{ha}$(z+ 的 Z= 曲. thah

\begin{equation*}\omega chah '\omega h\tag{6.55}\end{equation*}

式中 A 是缓变水域中线性波的波幅，它将由一阶近似方程确定. (2) 一阶近似解。( -i$\mu$) 的方程和边界条件为 的 -aZ t/> 1 =$\alpha$ 学 +JUdo 〉， -hO 〈 O dX d :r z = 0: 弘一 ?1=-MJt 呻川 /g

z =-h: 到 = .p o a 丝<I.r 一阶近似的速度势是非齐次线性常微分方程，而非齐次线性常微分方程的边值问题 不一定都有解，只有当右端函数满足可解性条件时才有解，可解性条件可由以下积分 公式导出.设在 (-11 .0) 中有任意两连续函数札 $\phi$ ，则有以下等式 z joh 「U 「 d ， d1[.p(I' - a 2 $\psi$) 一 $\psi$ ( rj\{' -a2 .p )Jdz = 1 .1 了(部 )-r( 甜') Idz = [哗'一份 'J I 气J 10 Lm二 Qz. J 令 .p= 仇 .${\rm ø}$= 仇，在水波零阶近似中有 .p o.. - a2 .p0 = 0 .因而有 [h.p" (的- a2.p1 ) dz = (.p" .p'1 - .pl .p;,) I$\upsilon$h il! Acha(z+h) 已知仇二监2 一一一一一，并用一阶近似的第一式替代职分式中的(旷-a 2 .p 1 )，将一阶

\[\omega ch\alpha h\]

近似的第二、第三式和线性水波的边界条件代人上式右边，得 fh (a 9? o 垫+仇主(忡。 ))dz=-l\{ 仇钱" + (w.po ), ) \}可悄)=\_. ha 翌h ,-T " dX . . " ${\rm e}$J.l: ' -. '" ,-- K 1 ' 1.' "......."H '''''''''U'l' '.::=V " " '= -$\tilde{n}$"" ${\rm e}$J.1 应用积分求导公式 z jo 川 dz=joh 叫 )dz-t$\alpha$耐 将式 (6. 56) 左端积分和右端最后一项合并，得 jo 卢 dz+ 拉叫 )= $\tilde{O}$ =。 由于仇=二豆主生豆豆土豆2 ，可得

\[\omega ch\alpha h\]

r 卢z=idfthd244$\omega$2$\omega$2 ch2ah (w .p \textasciitilde{} )= .0 = gA 2 缸， 再将等深度线性水波中群速度公式 C.=n g thah+ 斗孚7 代入得

\[L\omega \omega cn\alpha n\]

(6.56) \_d I$\iota$gA2l+ 乒\{皇生:\}=0 (6.57a > cJ.r L$\omega$J rlt 飞 $\omega$1 已知单位长度的波能 E=pgN/2. 因此缓变水深中水波能量的传播方程为 主 (c" : )十二(:)=。 由于.r =$\mu$.r .t - $\mu$t 是简单的线性变换，将式 (6.57a) 中慢变量恢复到普通变量后，有 d I E 、， rJ I E \textbackslash\{\} 1+ :: I:::\_ 1 = 0 (6.57c) d .r 飞民 $\omega$ J ' dt \textbackslash\{\} $\omega$/ 式中 E/$\omega$ 称为波作用量.式 (6.57a) 是缓变水域中波能量的演化方程，它的意义是披 作用量在群速度的轨迹上是守恒的.对于一般的平面行进水波在缓变地形 h=( $\mu$.r. $\mu$y) 中传播时，波数向量 k= $\alpha$ i+ 卢!j .群速度 c，， =c，，( $\mu$.1" .$\mu$ 川，这时波能传播方程为

J I E 飞 ，r1 1 E\textbackslash\{\} ,r1 IE\textbackslash\{\} \_ \textasciitilde{}一 I E 、， J I E 飞. 1 c.. - 1-卡 -:--IC... - I-卡 -1 - 1 = U !!Y. V. 1 C.. - 1 -$\leftarrow$一- 1 -一 1 = lJ J :r 飞 P$\omega$ 1 ' r1y \textbackslash\{\} -.."$\omega$ 1 ' r1 t \textbackslash\{\} $\omega$ 1 \textasciitilde{}币. \textbackslash\{\} -.$\omega$ 1 ' ."J, \textbackslash\{\}$\omega$ 1 \textasciitilde{} (6.57d) 给定缓变地形，将被能传播方程和波数连续方程(式 (6. 51) 或式 (6. 52 川， 式 (6.52b)) 联立，就可以计算在变深度浅水中波能传播.具体计算方法可参阅水波 动力学的专著.这里，可以简单地讨论水波进入浅水域的特性.当水波进入浅水域时， 被速和波群速将减小，为了满足波能传播方程，波能的陡度 (JE/$\theta$ :r) 将增加，这就是 通常见到的海浪进入歧滩时的涌浪现象.

\section*{6.5  线性浅水长波}\addcontentsline{toc}{section}{6.5 线性浅水长波}

水深和波 K 相比很小的水波称为浅水波 z 如果民披的波幅又比水深小得多，这种 水波称为线性浅水长波，例如 z 近海潮沙披属于浅水长波.线性浅水长波中.有以下 的数量关系 z

\begin{equation*}A >> h >> a\tag{6.58}\end{equation*}

根据浅水长波的特征，式 (6.58> 可以简化理想流体运动方程而得到解析解.为了分析 简明起见，仍假定波场是二维的，只有水平和垂直两个方向的运动，同时水深等于常数. 1.线性浅水长波的基本方程 二维理想流体运动的基本方程为 $\tilde{Ju}$ , r1 w L二+ 一二 = 0

\begin{gather*}
r1:r ' rJZ:\tag{6.59}\\
Du rJ u , Ju , rJ U 1\theta b
\end{gather*}

一一 = 一 +u 一 十 'U' 一=一 \textasciitilde{}"Y (6.60a) Dt rJ t ' - J:r ' - cJz: p J:r 旦旦 = d\_W + u d\_\_w + W \textasciitilde{}旦= - \_!\_ \textasciitilde{}主-\textasciitilde{} (6.60b) Dt r1 t' - r1:r ' - cJ z P r1 z n 利用浅水长波的基本性质，可以对 二 维理想流体运动的基本方程进行简化.令

\[:r = \lambda r. z - hz\]

.r .z 分别为流向和垂向的无量纲坐标.通过分析，可得浅水长波的以下性质: (1)流向速度在垂向近似不变 等深度线性水波的流向速度为式 (6.31al: u - 一旦旦;-:-cha ( Z: + h) cOs(a, r - wl )

\[sh(\alpha  h)\]

将 cha(z+h) 在 z 方向作泰勒展开 z 内 11 1 , , z 、 1 12$\pi$h/" z\textbackslash\{\}l cha(z +的 = ch ':;$\tilde{II}$ (1 + \textasciitilde{} ) = 1 +一|一". 11 + 7- 1 I +高阶小量A 飞 111 ., 2L${\rm A}$ 飞 h 1 J (6.6 1)

浅水波中 h/)'<<l; 在 -h$\tilde{z}$$\zeta$0 7.1<域中 O+z/的运 1 ，于是

\[Mz+h)=l+0(?)z\]

因此，在水深方向，流向速度可以近似等于常数(准确到 2 阶小量) ，即 u( .r， z) 句 u( .r) (6.62) (2) 垂向速度较流向速度小一量级 将连续性方程沿垂向积分，可得水域中的垂向速度〈线性水波理论中，波面可近 似为 z=O) :

\[Z Ju u-r a\]

一$\theta$ 且" ''BBEEFd 一一切 将.r=扯 ， z=hz 代人积分式，得 Z JU

\[\omega -H riltJ h-A\]

一w (6.63) 从上式可见，垂向速度较流向速度至少小一个量级，容易证明垂向加速度也较流向加 速度小一个量级. (3) 运动方程的简化 垂向运动方程中略去垂向质点加速度，得近似方程如下 z

\[-.l Jp = "\]

p Jz " 将上式沿 z 方向积分，在波面上 $\rho$= 队，得 $\rho$ = pg('\_ z) 十 $\rho$a 将压强代入流向运动方程，并略去非线性对流项，得 (6. 64) 笠 hg 一

\[h-h\]

(6. 65) (4) 水面边界条件的简化 水面边界条件为 w 一一 笠 hu + 当归 它的线性近似为 也一一 坐 h 由式 (6.63) 可知垂向速度 h-h 'n 一z 'd u--2$\upsilon$- 吨。

\["Ea--EE\]

d 一一切 于是水面边界条件简化为 句 -hLH 一一 笠归 (6.66)

习题 231 式 (6.65) 和式 (6.66) 是等深度线性长披的基本方程.将以上两式分别消去 u 和\textasciitilde{}，可得

\begin{gather*}
~-J!h ~ = 0 (6.67) d l" dx'\\
nu -
\end{gather*}

MAE 2 2

\[-z\]

吨。-吨。

\[LH g u-Z J-M\]

(6.68) 它们都是标准的双曲型方程，令 c 2 =gh ， 以上方程的一般解为(读者自己验证)

\begin{gather*}
~ = h[F(x-ct) - f(x+c t> J\\
u = c[F(x-ct) + f(x+a)J
\end{gather*}

F( 沪 ， f( 守)为有二阶导数的任意画数. 2. 线性浅水长波的性质 (1)线性浅水长波是非色散波 式 (6. 69) 和式 (6. 70) 表明，波面和流向速度都是以等速传播，其中函数 F(x-ct) 表示向右传播 z 函数 f(x+ct) 表示向左传播.波速 c= 品百 (6.71) 披速只和水深有关，它是非色散披，就是说，波形在传播过程中保持不变. (2) 波形由初始波面函数和流向速度确定 设初始波形 \textasciitilde{}(x ， O) = Z( 川，初始流向速度 u(x ， O) =U(x) ， 代入式 (6. 69) 和 式 (6.70) 得 (6.69) (6.70)

\begin{gather*}
Z(x) = h[F(x) - f(x)J\\
U(x) = c[F(x) + f(x)J
\end{gather*}

以上两式可解出函数 F(x) 和 f(x): 1 ,U(x) , Z(x) \textbackslash\{\} F(x) = : I 一一一十\_'.- , I !:飞 c n I (6. 72a) 1 ,U(x) Z(x) 飞 f(x) = \textasciitilde{} (一一一一\textasciitilde{}\textasciitilde{}' 1 (6.72b) L 飞 c h J (3) 缓变地形中线性浅水长波 只要将水面速度 u 用水深平均速度 u 代替，以上所得的等深度线性浅水长波结 果可以推广到缓变地形中线性浅水长波的情况，水深平均速度的定义是

\begin{gather*}
dbd=i| o uu , z , tdz\\
n J - h
\end{gather*}

(6. 73) 关于水波理论在海洋科学、造船和水利工程中的应用请参阅有关专著.

\section*{习题}\addcontentsline{toc}{section}{习题}

\noindent\textbf{6.1.} 求波长为 145m 的海洋波传播速度和波动周期，假定海洋深度是很大的-

\noindent\textbf{6.2.} 海洋波以 10m/s 的速度移动，试求这些波的波长和周期.

\noindent\textbf{6.3.} 在无限深液体表面，观察到表面上的浮漂 lmin 上升下降 15 次，试求波长

和传播速度.

\noindent\textbf{6.4.} 在无限深流体中，若置入两块无限大的铅垂平板，其垂直距离为 L ，试求当

两平板之间的流体产生驻波时，其波长 A 与 L 必须满足的关系，并求出其波动周期.

\noindent\textbf{6.5.} 在水深为 d 的水平底部(即 z=-d 处)用压力传感器记录到沿 r 方向传

播的周期行波的压力为 $\rho$ (t) ， 设 $\rho$ (t) 的最大脉动幅度(相对平衡态而言)为 H( 水柱 高) ，圆频率为 $\omega$ ，试确定所对应的自由面波动的圆频率和振幅.

\noindent\textbf{6.6.} 验证等深度线性浅水长波的解为式 (6.69) 和式 (6.70).

\noindent\textbf{6.7.} 设等深度线性浅水长波的初始波形\textasciitilde{}(.r , 0) = cos(.r) , u( .r, 0) = sin( 川，计

算该波形的传播.

\section*{思考题}\addcontentsline{toc}{section}{思考题}

S6.1 海浪冲向海岸时，浪高会越来越高，为什么? S6.2 在海底发生地震，为什么产生很高的海浪? S6.3 在水面下等速航行的潜艇，是再会在水面上出现水波? S6.4 铺设在海底的油管受到的海水作用力是恒定的，还是周期的?
